\scanpage{3}\mspage{\handwritten{2}}

On a aussi un foncteur évident
\[
T:(\mathrm{Ssimpl})\handwritten{\longrightarrow}(\mathrm{Cat})\ ,
\]
en associant à tout ensemble semi-simplicial $X$ la catégorie $T(X)=\handwritten{\Delta}_{/X}$ des simplexes sur $X$, dont l'ensemble des objets est la réunion disjointe des $X_n$ .... (c'est une catégorie fibrée sur la catégorie $\handwritten{\Delta}$ des simplexes types, à fibres les catégories discrètes définies par les $X_n$). Ceci posé, Quillen prouve que pour tout $X$, $ST(X)$ est isomorphe canoniquement à $X$ \unsure{dans} la catégorie homotopique construite avec $(\mathrm{Ssimpl})$, et que pour toute $C$, la catégorie $TS(C)$ est \inserted{canoniquement} ``homotopiquement équivalente à $C$'' \deleted{dans le sens} i.e. \inserted{canoniquement} isomorphe à $C$ dans la catégorie quotient de $(\mathrm{Cat})$ obtenue en inversant les foncteurs qui sont des homotopismes. \deleted{Il en résulte} Ces isomorphismes sont fonctoriels en $X$. Il en résulte formellement qu'un morphisme $f:X\handwritten{\unsure{-|\to}}Y$ dans $(\mathrm{Ssimpl})$ est un homotopisme si et seulement si il en est ainsi de $T(f):\unsure{T(X)}\handwritten{\to}T(Y)$, d'où des foncteurs $S':(\mathrm{Cat})'\handwritten{\longrightarrow}(\mathrm{Ssimpl})'$ et $T':(\mathrm{Ssimpl})'\handwritten{\to}\unsure{(\mathrm{Cat})'}$ entre les catégories ``homotopiques'', \inserted{construites avec $(\mathrm{Cat})$ resp. $(\mathrm{Ssimpl})$,} qui sont quasi-inverses l'un de l'autre.

De plus, Quillen construit un isomorphisme canonique et fonctoriel dans $(\mathrm{Cat})'$ entre $C$ et la catégorie opposée $C^o$, ou ce qui revient au même, un isomorphisme canonique et fonctoriel dans $(\mathrm{Ssimpl})'$ entre $S(C)$ et $S(C^o)$. \deleted{Cela lui permet alors, pour} La définition est telle que le foncteur induit sur les systèmes locaux sur $C$ tra\unsure{hsfôrme|hsforme|nsforme} le foncteur contravariant $F$ sur $C$, transformant toute flêche en flêche inversible, en le foncteur covariant (i.e. contravariant sur $C^o$) ayant mêmes valeurs sur les objets, et obtenu sur les flêches en remplaçant $F(u)$ par $F(u)^{-1}$; en d'autres termes, l'effet de l'homotopisme de Quillen sur les groupoïdes fondamentaux est l'isomorphisme évident entre les groupoïdes fondamentaux de $C$ et de $C^o$, \deleted{\unsure{qui est l'id en}} compte tenu que le deuxième est l'opposé du

% archive sheet number: 2
